%% notes-tex/multimodal/043-training-plan/sections/0-abstract.tex
%% SOURCE: results/overnight_2026_10_22_window_means.json (window means, seed 0), through the
%% design review's tables/t1-measured-rows.tex; budget totals from sections/5-grid.tex, whose
%% arithmetic is written out there from the sacct elapsed times of jobs 6136, 6137 and 6152.

\begin{abstract}
\noindent
Experiment 043 trains in two stages. Stage 1 pretrains one model over every dataset of the
043 query that has a head, with Ohya 2005 morphology read only as a probe of the frozen
model; it carries the arms that change the representation: sampler and weighting, the
where the interaction graphs act (a masked layer,
a readout gather, or nowhere), background-strain marks, the isolate route and the mark cap.
Stage 2 fine-tunes the stage-1 model on one target at a time, Kemmeren 2014 first, and
compares it with training from scratch at the same total optimizer steps and at the same
fine-tune steps. The case for the split is the overnight round of 2026-10-22: at seed 0 the
single joint objective read below the Kemmeren-alone run at matched steps for two of three
arms, and 0.208 to 0.230 on the Sameith 2015 doubles against 0.335 alone, with no
validation loss logged. Every arm is called by the decision rule of the design review (a
paired gain of at least 0.02 in 5 of 6 seeds, with a resampled interval excluding zero).
The grid runs smokes on GilaHyper, six owner-approved jobs on Delta, the six-seed stage-1
grid and the Kemmeren fine-tunes on Polaris, and the fine-tuning grid on Aurora once the
trainer runs on Intel XPU, for an estimated 148 Polaris and 117 Aurora node-hours.
\end{abstract}
